\documentclass[tikz,border=8pt]{standalone}


% общий стиль (подключается через TEXINPUTS)
% tikz-style.tex (общий стиль для всех TikZ-картинок)
\RequirePackage[T1,T2A]{fontenc}
\RequirePackage[utf8]{inputenc}

% Palatino-подобный шрифт (как в MathJax часто воспринимается «более вытянутым»)
\RequirePackage{txfonts}
% txfonts даёт предупреждение по кириллице (T2A/txr не определён).
% Оставляем txfonts для математики, а текстовый шрифт возвращаем на Computer Modern,
% чтобы не искать T2A/txr и убрать предупреждения.
\renewcommand{\rmdefault}{cmr}
\renewcommand{\sfdefault}{cmss}
\renewcommand{\ttdefault}{cmtt}

\RequirePackage{tikz}
\RequirePackage{xcolor}
\usetikzlibrary{calc}

% цвета
\definecolor{lightgreen}{RGB}{214,234,214}
\definecolor{darkgreen}{HTML}{0B7A0B}
\definecolor{darkred}{HTML}{DF0000}

\definecolor{midgreen}{HTML}{2E9B2E}
\definecolor{palegreen}{RGB}{235,246,235}
\definecolor{olivegreen}{HTML}{3E6B2F}

\definecolor{lightred}{RGB}{255,220,220}
\definecolor{midred}{HTML}{E04545}
\definecolor{darkred2}{HTML}{A80000}

\definecolor{lightblue}{RGB}{220,232,255}
\definecolor{midblue}{HTML}{2E5BD6}
\definecolor{darkblue}{HTML}{0B2E8A}

\definecolor{lightgray}{RGB}{240,240,240}
\definecolor{midgray}{RGB}{170,170,170}
\definecolor{darkgray}{RGB}{80,80,80}

\definecolor{vtxfill}{RGB}{86,193,255}
\definecolor{vennfill}{RGB}{220,232,255}


% единый набор стилей TikZ
\tikzset{
    lab/.style={
        font=\fontsize{24}{32}\selectfont
    }
}


\usepackage{tikz}
\usetikzlibrary{arrows.meta}

\begin{document}
\begin{tikzpicture}[x=1cm,y=1cm,line cap=round,line join=round]

% main number line
\draw[line width=1.0pt, <->, >=Stealth] (-7,0) -- (7,0);

% ticks and labels only at -2, 0, 4
\foreach \x in {-2,0,4} {
  \draw[line width=1.0pt] (\x,0.12) -- (\x,-0.12);
  \node[font=\Large] at (\x,-0.55) {$\x$};
}

% interval segment above the line: [-2, 4)
\draw[darkgreen, line width=2.0pt] (-2,0.55) -- (3.85,0.55);

% filled dot at -2 (included)
\fill[darkgreen] (-2,0.55) circle (3.6pt);

% open circle at 4 (excluded)
\draw[darkgreen, line width=2.0pt] (4,0.55) circle (3.6pt);

% red arrows pointing to endpoints
\draw[darkred, line width=1.2pt, -{Stealth[length=3.2mm]}] (-2.6,1.3) -- (-2.05,0.72);
\draw[darkred, line width=1.2pt, -{Stealth[length=3.2mm]}] ( 4.0,1.3) -- ( 4.00,0.72);

% explanatory text
\node[font=\large, align=center, text width=7.2cm, execute at begin node={\hyphenpenalty=10000\exhyphenpenalty=10000\relax}] at (-4.2,2.05) {закрашенный кружок показывает, что\\$-2$ \textbf{включено} во множество};
\node[font=\large, align=center, text width=7.2cm, execute at begin node={\hyphenpenalty=10000\exhyphenpenalty=10000\relax}] at ( 4.2,2.05) {незакрашенный кружок показывает, что\\$4$ \textbf{не включено} во множество};

\end{tikzpicture}
\end{document}